\documentclass[12pt,a4paper]{article}
\usepackage{cmap}
\usepackage[T2A]{fontenc}
\usepackage[utf8]{inputenc}
\usepackage[russian]{babel}
\usepackage{amsmath,amssymb,geometry,graphicx,float,hyperref,array,longtable}
\usepackage[labelfont=bf]{caption}
\geometry{left=20mm,right=20mm,top=20mm,bottom=20mm}
\hypersetup{hidelinks,pdftitle={Вопросы и ответы к защите лабораторной 1 Вариант 11},
pdfauthor={Иванов Никита Михайлович}}
\captionsetup{skip=5pt}
\setlength{\emergencystretch}{2em}
\setlength{\intextsep}{8pt}
\setcounter{tocdepth}{1}
\clubpenalty=10000
\widowpenalty=10000
\raggedbottom
\newcommand{\norm}[1]{\left\lVert#1\right\rVert_2}
\newcounter{question}
\newcommand{\question}[1]{\refstepcounter{question}\subsection*{\thequestion.\ #1}}
\title{Контрольные вопросы и ответы\\Лабораторная работа №1\\Статическая оптимизация\\Вариант 11}
\author{Иванов Никита Михайлович\\Группа R3436}
\date{5 ноября 2026 г.}
\begin{document}
\maketitle
\newpage
\tableofcontents
\newpage

\section{Постановка задачи и основные понятия}
Материал организован от определений к доказательствам и численным методам.
В ответах используются данные варианта 11:
\[
J(x,u)=8x^2+6u^2+4xu+x+3u-10,\qquad c(x,u)=3x-u^2+1.
\]
Вектор переменных обозначается $z=(x,u)^{\mathsf T}$.
Первые предложения ответов можно использовать как начало устного объяснения;
формулы и последующие рассуждения раскрывают его обоснование.

\question{Какова цель работы и какие задачи в ней решаются}
Нужно найти минимум одного и того же критерия в трёх постановках:
на всей плоскости, на кривой $c=0$ и в области $c\leq0$.
Затем требуется найти безусловный минимум численно методом
Ньютона--Рафсона и градиентным методом при двух постоянных шагах.

В аналитической части важны не только координаты кандидатов, но и
доказательство того, что они дают глобальные минимумы. В численной части
сравниваются направления и величины шагов, характер сходимости и число
итераций при одной начальной точке и одинаковом условии остановки.

\question{Что такое критерий качества и ограничение}
Критерий $J$ сопоставляет каждой паре $(x,u)$ число, по которому
сравниваются решения. Здесь задача минимизации: меньшему $J$ соответствует
лучший вариант по выбранной математической оценке. Конкретная физическая
интерпретация этого полинома в задании не задана.

Ограничение определяет допустимость решения. При $c=0$ разрешены только
точки параболы $x=(u^2-1)/3$; при $c\leq0$ разрешены точки
$x\leq(u^2-1)/3$. Минимизируется $J$, а у $c$ проверяется выполнение
равенства или неравенства. Условный минимум выбирается среди допустимых точек.

\question{Почему задача называется статической оптимизацией}
Здесь выбираются значения переменных $x$ и $u$, а не их функции времени.
В постановке нет дифференциального уравнения движения объекта и нет
интегрального критерия по времени. Критерий является обычной функцией
конечного числа переменных.

Номер итерации $k$ обозначает шаг вычислительного алгоритма. Он не является
физическим временем движения системы. Даже если в других задачах $x$
обозначает состояние, а $u$ управление, в этой работе они выступают
переменными статической задачи.

\question{Что такое градиент и как он вычисляется}
Градиент --- столбец первых частных производных:
\[
\nabla J(z)=
\begin{pmatrix}J_x\\J_u\end{pmatrix}
=\begin{pmatrix}16x+4u+1\\4x+12u+3\end{pmatrix}.
\]
Производная $J_x$ вычисляется при фиксированном $u$, а $J_u$ ---
при фиксированном $x$. Например, производная $4xu$ по $x$ равна $4u$.
В точке $z_0=(0,0)^{\mathsf T}$ градиент равен $(1,3)^{\mathsf T}$.

Для малого смещения $h$ справедливо
$J(z+h)-J(z)=\nabla J(z)^{\mathsf T}h+o(\norm h)$.
Поэтому градиент описывает изменение критерия в первом приближении.
В евклидовой геометрии он направлен в сторону самого быстрого локального роста.

\question{Что такое стационарность и достаточно ли её для минимума}
Стационарность означает $\nabla J(z^*)=0$. Для дифференцируемой функции
во внутренней точке области определения это необходимое условие
безусловного локального минимума.

Оно не является достаточным: у функций $x^2+u^2$, $-x^2-u^2$ и
$x^2-u^2$ градиент в начале координат равен нулю, но там соответственно
минимум, максимум и седло. Чтобы установить тип стационарной точки,
исследуют вторые производные или непосредственно изменение функции.
При ограничениях условный минимум может иметь ненулевой градиент $J$.

\question{Что показывает матрица Гессе}
Гессиан состоит из вторых частных производных:
\[
H_J(z)=
\begin{pmatrix}J_{xx}&J_{xu}\\J_{ux}&J_{uu}\end{pmatrix}
=\begin{pmatrix}16&4\\4&12\end{pmatrix}.
\]
Он описывает кривизну функции по различным направлениям. Для функции
класса $C^2$ смешанные производные совпадают, поэтому гессиан симметричен.
Разложение около стационарной точки имеет вид
\[
J(z^*+h)=J(z^*)+\frac12h^{\mathsf T}H_J(z^*)h+o(\norm h^2).
\]
Положительная определённость $H_J(z^*)$ обеспечивает строгий локальный
минимум. Для нашей квадратичной функции остаточного члена нет:
равенство с квадратичной формой выполняется точно.

\question{Как точно формулируется критерий Сильвестра}
Вещественная симметричная матрица $A\in\mathbb R^{n\times n}$
положительно определена тогда и только тогда, когда все её
угловые главные миноры положительны:
\[
\Delta_k=\det A_{1:k,1:k}>0,\qquad k=1,\ldots,n.
\]
Здесь берётся подматрица из первых $k$ строк и первых $k$ столбцов.
Для матрицы второго порядка условия имеют вид
$a_{11}>0$, $\det A>0$.
В работе $\Delta_1=16$, $\Delta_2=16\cdot12-4^2=176$.

Положительная определённость означает $h^{\mathsf T}Ah>0$ при $h\ne0$.
Если один из указанных миноров нулевой, строгий критерий не выполнен;
это само по себе не доказывает отсутствие минимума. Например,
$x^4+u^4$ имеет строгий минимум при нулевом гессиане в начале координат.

\question{Чем отличаются локальный минимум, глобальный минимум и выпуклость}
Локальный минимум сравнивает значение функции только с допустимыми
точками некоторой окрестности. Глобальный минимум сравнивает его
со всеми допустимыми точками:
\[
J(z^*)\leq J(z)\quad\text{для любого допустимого }z.
\]
Строгая выпуклость на выпуклой области означает, что при $z\ne y$ и
$0<\theta<1$
\[
J(\theta z+(1-\theta)y)
<\theta J(z)+(1-\theta)J(y).
\]
У строго выпуклой функции на выпуклом допустимом множестве может быть
не более одной точки минимума. Существование минимума нужно обосновать
дополнительно. В нашей задаче $H_J$ положительно определён на всей плоскости,
поэтому $J$ строго выпукла. Для невыпуклого ограничения этот факт
не заменяет отдельного доказательства условного глобального минимума.

\section{Минимум без ограничений}

\question{Как получить координаты безусловного минимума}
Приравниваем обе компоненты градиента к нулю:
\[
\begin{cases}16x+4u+1=0,\\4x+12u+3=0.\end{cases}
\]
Умножим второе уравнение на 4 и вычтем из него первое:
$44u+11=0$, откуда $u=-1/4$.
Подстановка в первое уравнение даёт $16x=0$, то есть $x=0$.
Таким образом,
\[
z_{\mathrm{free}}^*=(0,-1/4)^{\mathsf T},\qquad
J_{\mathrm{free}}^*=-10.375.
\]
Это пока стационарная точка; её минимальность следует из положительной
определённости гессиана.

\question{Как строго доказать глобальность и единственность этого минимума}
Для квадратичного критерия и найденной стационарной точки точно
выполняется
\[
J(z)-J(z_{\mathrm{free}}^*)
=\frac12(z-z_{\mathrm{free}}^*)^{\mathsf T}
H_J(z-z_{\mathrm{free}}^*).
\]
Поскольку $H_J\succ0$, правая часть положительна при любой другой точке.
Значит, меньшего или равного значения в другой точке быть не может.

То же видно после выделения квадратов:
\[
J(x,u)=-\frac{83}{8}
+8\left(x+\frac{u+1/4}{4}\right)^2
+\frac{11}{2}\left(u+\frac14\right)^2.
\]
Оба квадрата одновременно обращаются в нуль только при $(0,-1/4)$.
Это непосредственное доказательство единственного глобального минимума.

\question{Почему минимум существует и допустимо ли отрицательное значение критерия}
Отрицательное значение допустимо: в задании $J$ является математическим
полиномом с постоянным слагаемым $-10$. Оно не обязано быть физической
энергией. Добавление константы к критерию меняет минимальное значение,
но не меняет градиент и положение минимума.

Существование связано с ростом функции на бесконечности.
Запишем $J(z)=\frac12z^{\mathsf T}H_Jz+b^{\mathsf T}z-10$,
где $b=(1,3)^{\mathsf T}$. Если $m=\lambda_{\min}(H_J)>0$, то
\[
J(z)\geq\frac m2\norm z^2-\norm b\,\norm z-10
\longrightarrow+\infty\quad\text{при }\norm z\to\infty.
\]
Непрерывная функция с таким свойством достигает минимума.
Строгая выпуклость обеспечивает его единственность на всей плоскости.

\question{Может ли добавление ограничения уменьшить минимальное значение}
Добавление ограничения сужает выбор, поэтому минимальное значение
уменьшиться не может. Для наших трёх множеств
\[
\{c=0\}\subseteq\{c\leq0\}\subseteq\mathbb R^2
\]
следует порядок
\[
J_{\mathrm{free}}^*\leq J_{\leq}^*\leq J_{=}^*.
\]
В работе $J_{\leq}^*=J_{=}^*=-9.604836264941$, а безусловное значение
равно $-10.375$. Ограничение увеличило минимум на $0.770163735059$.
Увеличение означает ухудшение значения минимизируемого показателя;
оно остаётся лучшим среди допустимых решений.

\section{Ограничение в виде равенства и метод Лагранжа}

\question{Что такое лагранжиан и чем он отличается от критерия}
Критерий $J$ оценивает качество исходного решения.
Лагранжиан объединяет критерий и ограничение для записи условий
оптимальности:
\[
L(x,u,\lambda)=J(x,u)+\lambda c(x,u).
\]
Это скалярная функция с дополнительной переменной $\lambda$,
называемой множителем Лагранжа. Градиент и гессиан сами по себе
не являются лагранжианом: это производные соответствующей функции.

На допустимой кривой $c=0$ значение $L$ совпадает с $J$ при любом
$\lambda$, но их производные по $x,u$ обычно различаются.
При решении безусловной задачи множитель не нужен.

\question{Откуда берутся условия Лагранжа}
В регулярной точке кривой $c=0$ допустимое касательное направление $h$
удовлетворяет $\nabla c(z^*)^{\mathsf T}h=0$.
В условном минимуме производная критерия вдоль касательной равна нулю:
$\nabla J(z^*)^{\mathsf T}h=0$.
Поэтому $\nabla J$ и $\nabla c$ перпендикулярны одной касательной
и пропорциональны:
\[
\nabla J(z^*)+\lambda^*\nabla c(z^*)=0.
\]
Вместе с допустимостью это даёт условия
\[
L_x=0,\qquad L_u=0,\qquad L_\lambda=c=0.
\]
Необходимость предполагает $\nabla c(z^*)\ne0$.
У нас $\nabla c=(3,-2u)^{\mathsf T}$, поэтому регулярность выполняется
во всех точках границы.

\question{Почему множитель равенства имеет любой знак и по каким переменным берут гессиан}
У равенства нет выбранной допустимой стороны: движение должно оставаться
на кривой. Поэтому множитель $\lambda$ может иметь любой знак.
Равенство $L_\lambda=0$ возвращает само ограничение.

При проверке кривизны по переменным решения используется
\[
H_L=\nabla_{x,u}^2L
=\begin{pmatrix}16&4\\4&12-2\lambda\end{pmatrix}.
\]
Это матрица второго порядка по $x,u$ при фиксированном $\lambda$.
Полный гессиан по трём переменным $(x,u,\lambda)$ для этой проверки
не нужен. Метод Лагранжа не требует совместного минимума $L$ по всем
трём переменным: по $\lambda$ функция линейна.

Для локального условного минимума достаточно положительности
$h^{\mathsf T}H_Lh$ на ненулевых касательных направлениях.
Положительная определённость $H_L$ на всей плоскости является
более сильным достаточным условием.

\question{Как получить кубическое уравнение для условного минимума}
Необходимые условия имеют вид
\[
\begin{cases}
16x+4u+1+3\lambda=0,\\
4x+12u+3-2\lambda u=0,\\
3x-u^2+1=0.
\end{cases}
\]
Из третьего уравнения $x=(u^2-1)/3$, а из первого
\[
\lambda=-\frac{16u^2+12u-13}{9}.
\]
Подставляя оба выражения во второе уравнение, получаем
\[
\frac43(u^2-1)+12u+3
+\frac{2u}{9}(16u^2+12u-13)=0.
\]
После умножения на 9 и приведения подобных слагаемых остаётся
\[
32u^3+36u^2+82u+15=0.
\]
Это уравнение определяет координату $u$ стационарной точки на ограничении.

\question{Что означает функция \(\varphi(u)\) и как она получается}
Функция $\varphi$ является тем же критерием $J$, вычисленным только
на допустимой параболе. Каждое вещественное $u$ задаёт ровно одну
допустимую точку $\bigl((u^2-1)/3,u\bigr)$.
Поэтому параметризация охватывает всю кривую без потери решений.

Подстановка даёт
\begin{align*}
\varphi(u)&=8\left(\frac{u^2-1}{3}\right)^2+6u^2
+\frac{4u(u^2-1)}3+\frac{u^2-1}3+3u-10\\
&=\frac{8u^4+12u^3+41u^2+15u-85}{9}.
\end{align*}
Так задача двух переменных с равенством превращается в безусловную
одномерную задачу $\varphi(u)\to\min$. Уравнение $\varphi'(u)=0$
совпадает с кубическим уравнением, полученным через лагранжиан.

\question{Почему у кубического уравнения ровно один действительный корень}
Для функции $\varphi$
\[
\varphi''(u)
=\frac{96u^2+72u+82}{9}
=\frac{32}{3}\left(u+\frac38\right)^2+\frac{137}{18}>0.
\]
Следовательно, $\varphi'$ строго возрастает.
Её старший член равен $(32/9)u^3$, поэтому при $u\to-\infty$
производная стремится к $-\infty$, а при $u\to+\infty$ --- к $+\infty$.
Непрерывность обеспечивает хотя бы один нуль, а строгое возрастание
исключает второй.

Кроме того, $\varphi(u)\to+\infty$ при $|u|\to\infty$,
так как коэффициент при $u^4$ положителен.
Единственная стационарная точка задаёт единственный глобальный
минимум на всей параболе.

\question{Каковы условное решение и его основные проверки}
Численные приближения равны
\[
u^*\approx-0.196978642820,\quad
x^*\approx-0.320399804758,\quad
\lambda^*\approx1.638103815800,
\]
\[
J_{=}^*\approx-9.604836264941.
\]
Нужно проверить допустимость $3x^*-(u^*)^2+1=0$ и стационарность
$\nabla J(z^*)+\lambda^*\nabla c(z^*)=0$.
Для округлённых координат равенства выполняются с погрешностью округления.
В точке решения $\varphi''(u^*)\approx7.949154883>0$.

При этом сам градиент $J$ не равен нулю:
\[
\nabla J(z^*)\approx(-4.91431145,-0.64534293)^{\mathsf T}.
\]
Это нормально для условного минимума: его компенсирует
$\lambda^*\nabla c(z^*)$. Исчезновение свободного градиента
для такой точки не требуется.

\section{Неравенство и условия Куна--Таккера}

\question{Какие четыре группы условий входят в условия Куна--Таккера}
Для задачи $\min J(z)$ при $c(z)\leq0$ и лагранжиана
$L=J+\lambda c$ условия имеют вид
\[
\begin{aligned}
\nabla J+\lambda\nabla c&=0
&&\text{(стационарность)},\\
c&\leq0&&\text{(допустимость)},\\
\lambda&\geq0&&\text{(знак множителя)},\\
\lambda c&=0&&\text{(дополняющая нежёсткость)}.
\end{aligned}
\]
Они описывают связь критерия с ограничением, допустимость решения
и влияние активной границы. Знак множителя относится именно к записи
$c\leq0$ и $L=J+\lambda c$.

Для гладкой задачи это необходимые условия при выполнении условий
регулярности. При одном активном неравенстве достаточно
$\nabla c(z^*)\ne0$, что в нашем варианте выполнено.
Достаточность для глобального минимума в общем случае из этих условий
не следует; она доказывается отдельно.

\question{Что такое активное ограничение и может ли его множитель быть нулевым}
Ограничение активно, если $c(z^*)=0$, то есть точка лежит на границе.
Если $c(z^*)<0$, ограничение неактивно, и из $\lambda c=0$
следует $\lambda=0$. Тогда стационарность превращается в $\nabla J=0$.

Для активного ограничения множитель может быть положительным или нулевым.
Положительный множитель обязательно означает активность, но обратное
утверждение неверно. Например, у задачи
$\min_x x^2$ при $-x\leq0$ минимум $x^*=0$ лежит на границе,
а условие $2x-\lambda=0$ даёт $\lambda^*=0$.
Таким образом, нельзя заменять условие $\lambda\geq0$ строгим
неравенством для всех активных ограничений.

\question{Выпукла ли допустимая область в нашем варианте}
Допустимое множество имеет вид
\[
D=\left\{(x,u):x\leq\frac{u^2-1}{3}\right\}.
\]
Оно невыпукло. Точки $A=(0,1)$ и $B=(0,-1)$ допустимы:
в обеих $c=0$. Их середина $(0,0)$ недопустима, поскольку $c(0,0)=1>0$.
По определению выпуклого множества все отрезки между его точками
должны оставаться внутри; здесь это нарушается.

Сам критерий $J$ строго выпуклый, но это не превращает задачу
с невыпуклым множеством в выпуклую. Поэтому применение только
общей теоремы о достаточности ККТ для выпуклой задачи здесь
не было бы корректным.

\question{Как выбрать между неактивной и активной ветвями}
Сначала рассматриваем $\lambda=0$. Тогда получается безусловная
точка $z_{\mathrm{free}}^*=(0,-1/4)^{\mathsf T}$, но
\[
c(z_{\mathrm{free}}^*)=1-\frac1{16}=\frac{15}{16}>0.
\]
Она недопустима, поэтому эта ветвь не даёт решения задачи.

Далее рассматриваем $c=0$. Получается ранее найденная точка на параболе
с $\lambda^*\approx1.638103815800>0$. Она удовлетворяет всем
условиям Куна--Таккера.
В результате минимумы для равенства и неравенства совпали.
Это свойство данной задачи; в общем случае безусловный минимум
может находиться внутри допустимой области, и тогда результаты
для равенства и неравенства будут различаться.

\question{Как доказать глобальный минимум при невыпуклом ограничении}
Зафиксируем найденный $\lambda^*$ и рассмотрим $L(z,\lambda^*)$
как функцию $z=(x,u)^{\mathsf T}$. Её гессиан равен
\[
H_L=\begin{pmatrix}16&4\\4&12-2\lambda^*\end{pmatrix},
\qquad
\det H_L=176-32\lambda^*\approx123.580677894>0.
\]
Первый угловой минор равен 16, поэтому $H_L\succ0$.
Из стационарности следует, что $z^*$ --- единственный глобальный
минимум $L(z,\lambda^*)$ на всей плоскости.

Для любой допустимой точки $\lambda^*c(z)\leq0$, поэтому
\[
J(z)\geq L(z,\lambda^*)\geq L(z^*,\lambda^*)=J(z^*).
\]
Последнее равенство использует $c(z^*)=0$.
Для $z\ne z^*$ второе неравенство строгое, что доказывает
единственность условного глобального минимума.
В этом доказательстве выпуклость допустимой области не требуется.

\question{Каков смысл множителя \(\lambda^*\)}
Геометрически множитель обеспечивает равенство
$\nabla J=-\lambda^*\nabla c$ на активной границе.
Он показывает, как граница компенсирует свободное направление
уменьшения критерия.

Есть также локальная интерпретация чувствительности. Заменим
ограничение на $c(z)\leq\delta$ и обозначим оптимальное значение
через $V(\delta)$. При гладком продолжении того же регулярного
оптимального решения
\[
V'(0)=-\lambda^*,\qquad
V(\delta)\approx V(0)-1.638103816\,\delta.
\]
Положительное $\delta$ ослабляет ограничение, расширяет допустимый выбор
и локально уменьшает оптимальное значение. Это оценка для малых
изменений на устойчивой ветви решения, а не точная формула
для произвольного изменения ограничения.

\section{Метод Ньютона--Рафсона и остановка поиска}

\question{Откуда берётся формула Ньютона для одной переменной}
При оптимизации метод Ньютона решает уравнение $J'(x)=0$.
Линеаризуем производную около текущей точки:
\[
J'(x_k+\Delta x)\approx J'(x_k)+J''(x_k)\Delta x.
\]
Требуем равенства нулю этого линейного приближения. При
$J''(x_k)\ne0$ получаем
\[
\Delta x=-\frac{J'(x_k)}{J''(x_k)},\qquad
x_{k+1}=x_k+\Delta x.
\]
Шаг $\Delta x$ обычно не равен нулю. Нулю приравнивается приближение
производной в новой точке. Сам шаг становится нулевым,
когда текущая производная уже равна нулю.
Для неквадратичной функции новая истинная производная обычно
не обнуляется точно за один шаг.

\question{Как получается многомерная формула Ньютона}
Градиент линеаризуется по формуле
\[
\nabla J(z_k+s)\approx g_k+H_ks,\qquad
g_k=\nabla J(z_k),\quad H_k=\nabla^2J(z_k).
\]
Условие нулевого приближения даёт систему $H_ks=-g_k$.
Если $H_k$ обратим,
\[
z_{k+1}=z_k-H_k^{-1}g_k.
\]
Это также стационарная точка квадратичной модели
$J(z_k)+g_k^{\mathsf T}s+\frac12s^{\mathsf T}H_ks$.
Если $H_k\succ0$, стационарная точка этой модели является её минимумом.
Градиент учитывает наклон функции, а гессиан изменяет шаг
в соответствии с кривизной по различным направлениям.

\question{Как вручную проверить единственный шаг в работе}
В начальной точке
\[
z_0=(0,0)^{\mathsf T},\qquad
g_0=(1,3)^{\mathsf T},\qquad
H^{-1}=\frac1{176}\begin{pmatrix}12&-4\\-4&16\end{pmatrix}.
\]
Произведение равно
\[
H^{-1}g_0
=\frac1{176}\begin{pmatrix}12-12\\-4+48\end{pmatrix}
=\begin{pmatrix}0\\1/4\end{pmatrix}.
\]
Следовательно, $z_1=(0,-1/4)^{\mathsf T}$.
При подстановке в градиент получаем
$(16\cdot0+4(-1/4)+1,\ 4\cdot0+12(-1/4)+3)^{\mathsf T}=0$.
Значение критерия изменилось с $J_0=-10$ до $J_1=-10.375$.
Нулевой градиент устанавливает стационарность, а положительно
определённый гессиан устанавливает минимум.

\question{Почему для квадратичной функции хватает одного шага из любой точки}
У критерия $J(z)=\frac12z^{\mathsf T}Hz+b^{\mathsf T}z-10$
градиент равен $Hz+b$, а гессиан $H$ постоянен.
Поэтому
\[
z_{k+1}=z_k-H^{-1}(Hz_k+b)=-H^{-1}b=z^*.
\]
Это точное тождество: квадратичная модель совпадает с исходной функцией.
В точной арифметике из любой начальной точки достаточно одного
полного шага при обратимом $H$.

Постоянство и обратимость гессиана объясняют один шаг к стационарной
точке. Положительная определённость объясняет её минимальность.
Например, у $x^2-u^2$ гессиан также постоянен и обратим,
но его стационарная точка является седлом.

\question{Что означает норма градиента и почему по ней останавливаются}
Евклидова норма градиента равна
\[
\norm{\nabla J}=\sqrt{J_x^2+J_u^2}.
\]
Индекс 2 обозначает вид нормы, а не возведение всего выражения в квадрат.
Условие $\norm{\nabla J(z_k)}\leq10^{-6}$ означает, что обе компоненты
градиента малы в совокупности. Число $10^{-6}$ --- выбранный порог,
а не универсальная константа метода. Для Ньютона в работе норма равна нулю.

Для нашего положительно определённого квадратичного критерия
$g_k=H(z_k-z^*)$. Если $m=\lambda_{\min}(H)$, то
\[
\norm{z_k-z^*}\leq\frac{\norm{g_k}}m,\qquad
0\leq J(z_k)-J^*\leq\frac{\norm{g_k}^2}{2m}.
\]
При пороге $10^{-6}$ эти верхние оценки равны примерно
$1.05\cdot10^{-7}$ для расстояния и $5.25\cdot10^{-14}$ для ошибки
критерия. Такие оценки используют известную положительную кривизну;
малый градиент сам по себе не доказывает минимум произвольной функции.

\question{Всегда ли метод Ньютона сходится к минимуму}
Для произвольной функции такой гарантии нет.
Гессиан может оказаться необратимым; стационарная точка может быть
максимумом или седлом; из далёкой начальной точки полный шаг может
увеличить критерий или привести к расходимости.

Если гессиан локально липшицев в окрестности решения и в решении
невырожден, метод имеет локальную квадратичную сходимость из достаточно
близких начальных точек. Для минимизации обычно обеспечивают положительную
кривизну модели и при необходимости уменьшают шаг.
В нашей задаче гессиан всюду постоянен и положительно определён,
поэтому эти затруднения отсутствуют и действует точный результат
об одном шаге.

\section{Градиентный поиск и выбор шага}

\question{Почему направление минус градиент даёт наискорейшее уменьшение}
Производная по единичному направлению $d$ равна
$D_dJ=g^{\mathsf T}d$, где $g=\nabla J(z)$.
По неравенству Коши--Буняковского
\[
g^{\mathsf T}d\geq-\norm g\,\norm d=-\norm g.
\]
При $g\ne0$ нижняя граница достигается для
$d=-g/\norm g$. Поэтому минус градиент задаёт наиболее быстрое
локальное уменьшение среди направлений одинаковой длины
в евклидовой норме.

Итерация имеет вид $z_{k+1}=z_k-\gamma_k g_k$.
Параметр $\gamma_k$ масштабирует направление, а длина фактического
перемещения равна $\gamma_k\norm{g_k}$.
Локальная крутизна не гарантирует уменьшения при произвольно
большом конечном шаге: величину шага нужно выбирать отдельно.

\question{Как записать итерации и проверить первые два шага}
Для варианта 11
\[
\begin{cases}
x_{k+1}=x_k-\gamma(16x_k+4u_k+1),\\
u_{k+1}=u_k-\gamma(4x_k+12u_k+3).
\end{cases}
\]
При $\gamma=0.05$ из $z_0=(0,0)^{\mathsf T}$ и
$g_0=(1,3)^{\mathsf T}$ получаем
$z_1=(-0.05,-0.15)^{\mathsf T}$.
В этой точке новый градиент равен $g_1=(-0.4,1)^{\mathsf T}$.
Поэтому
\[
z_2=(-0.05,-0.15)^{\mathsf T}
-0.05(-0.4,1)^{\mathsf T}=(-0.03,-0.20)^{\mathsf T}.
\]
При $\gamma=0.09$ первые точки равны
$z_1=(-0.09,-0.27)^{\mathsf T}$ и
$z_2=(0.0468,-0.216)^{\mathsf T}$.
Градиент пересчитывается в каждой новой точке. Одновременное обновление
обеих координат использует одну и ту же старую точку $z_k$.

\question{Как выводится условие сходимости постоянного шага}
Пусть $z^*$ --- безусловный минимум и $e_k=z_k-z^*$.
Поскольку $Hz^*+b=0$, имеем
\[
e_{k+1}=(I-\gamma H)e_k.
\]
Симметричная матрица $H$ имеет ортонормированный базис собственных
векторов. В каждой собственной координате ошибка умножается
на $r_i=1-\gamma\mu_i$, где $\mu_i$ --- собственное значение $H$.
Для сходимости из любой начальной точки необходимо и достаточно
$|r_i|<1$ для всех $i$. При $H\succ0$ это равносильно
\[
0<\gamma<\frac2{\mu_{\max}}.
\]
У нас характеристическое уравнение $\mu^2-28\mu+176=0$, поэтому
\[
\mu_{\min}=14-2\sqrt5\approx9.527864045,\qquad
\mu_{\max}=14+2\sqrt5\approx18.472135955,
\]
\[
0<\gamma<0.108271182330.
\]
При $\gamma=0$ движения нет. На верхней границе один множитель
равен $-1$ и соответствующая составляющая ошибки не затухает.
При ещё большем шаге она возрастает по модулю.

\question{Что означает апериодическая и колебательная сходимость}
Характер изменения определяется знаками множителей $r_i$.
Если $0<r_i<1$, соответствующая собственная составляющая ошибки
сохраняет знак и убывает по модулю. Это апериодический режим.
Если $-1<r_i<0$, составляющая меняет знак на каждом шаге,
а её модуль уменьшается: появляются затухающие колебания.

Для всех собственных составляющих апериодичность гарантируется при
$0<\gamma<1/\mu_{\max}$; при равенстве верхней границе самая быстрая
составляющая исчезает за один шаг.
В интервале $1/\mu_{\max}<\gamma<2/\mu_{\max}$
хотя бы один множитель отрицателен.
Это утверждения о собственных координатах ошибки.
Обычная координата $x_k$ или $u_k$ может быть суммой нескольких
составляющих, поэтому её монотонность отдельно не гарантируется.

\question{Почему выбраны шаги 0.05 и 0.09 и почему различается число итераций}
Порог смены знака равен $1/\mu_{\max}\approx0.054135591165$.
При $\gamma=0.05$ множители равны
\[
r_1\approx0.523606798,\qquad r_2\approx0.076393202;
\]
оба положительны. При $\gamma=0.09$ получаем
\[
r_1\approx0.142492236,\qquad r_2\approx-0.662492236;
\]
это устойчивый колебательный режим.

Спектральный радиус матрицы итерации равен
$\rho=\max_i|r_i|$. Для первого шага он примерно $0.5236$,
для второго --- $0.6625$. Поэтому в данном сравнении первый режим
асимптотически быстрее: он сильнее сжимает самую медленную составляющую.
Из одинаковой точки при пороге $10^{-6}$ получены 23 и 36 итераций.
Начальное состояние имеет номер $k=0$; число итераций равно числу
переходов к последнему состоянию.

\question{Почему координаты колеблются, а критерий всё равно убывает}
Для квадратичного критерия и $g=\nabla J(z)$ точно выполняется
\[
J(z-\gamma g)-J(z)
=-\gamma\norm g^2+\frac{\gamma^2}{2}g^{\mathsf T}Hg.
\]
Так как $g^{\mathsf T}Hg\leq\mu_{\max}\norm g^2$, получаем
\[
J(z-\gamma g)-J(z)
\leq-\gamma\left(1-\frac{\gamma\mu_{\max}}2\right)\norm g^2.
\]
При $0<\gamma<2/\mu_{\max}$ правая часть отрицательна для $g\ne0$.
Следовательно, критерий строго уменьшается до достижения минимума.

При этом переход через оптимум по одной координате возможен.
Знак ошибки и её вклад в значение функции --- разные характеристики.
В собственных координатах ошибка критерия является суммой положительных
слагаемых $\mu_i e_i^2/2$, поэтому смена знака $e_i$ не меняет знак
этого вклада.

\question{Чем постоянный шаг отличается от точного наискорейшего спуска}
В работе направление равно минус градиенту, а $\gamma$ заранее
фиксируется отдельно для каждого опыта. Это градиентный спуск
с постоянным шагом.

При точном выборе шага на каждой итерации решают одномерную задачу
\[
\gamma_k=\arg\min_{\gamma\geq0}J(z_k-\gamma g_k).
\]
Для положительно определённого квадратичного критерия и $g_k\ne0$
дифференцирование по $\gamma$ даёт
\[
-g_k^{\mathsf T}g_k+\gamma g_k^{\mathsf T}Hg_k=0,\qquad
\gamma_k=\frac{g_k^{\mathsf T}g_k}{g_k^{\mathsf T}Hg_k}.
\]
Из начальной точки первый точный шаг равен $10/148=5/74\approx0.067567568$.
Он отличается от обоих шагов работы.
Если на защите уточнят терминологию, следует явно назвать выбранный
вариант: поиск в направлении минус градиента с постоянными шагами
$0.05$ и $0.09$. Формула точного выбора шага описывает другой вариант
алгоритма.

\question{Можно ли выбрать постоянный шаг лучше использованных}
Для положительно определённой квадратичной функции постоянный шаг,
минимизирующий наихудший коэффициент сжатия ошибки, равен
\[
\gamma_{\mathrm{opt}}=\frac2{\mu_{\min}+\mu_{\max}}.
\]
Он получается из уравнивания модулей крайних множителей:
$1-\gamma\mu_{\min}=-(1-\gamma\mu_{\max})$.
В нашей задаче
\[
\gamma_{\mathrm{opt}}=\frac1{14}\approx0.071428571,\qquad
\rho_{\mathrm{opt}}=\frac{\mu_{\max}-\mu_{\min}}
{\mu_{\max}+\mu_{\min}}=\frac{\sqrt5}{7}\approx0.319438282.
\]
Это оптимальность по наихудшему сжатию среди всех начальных ошибок,
а не решение одномерного поиска шага на каждой итерации.
Заданные опыты предназначены для сравнения двух характеров сходимости;
они не утверждают оптимальности выбранных постоянных параметров.

\question{Почему численные методы сходятся в точку, нарушающую ограничение}
В численной части работы минимизируется $J_1=J$ без ограничений.
Поэтому алгоритмы Ньютона и градиентного поиска должны прийти
к $z_{\mathrm{free}}^*=(0,-1/4)^{\mathsf T}$.
Нарушение $c\leq0$ в этой точке не является ошибкой данной части:
она решает другую постановку, чем задача с неравенством.

Если бы требовалось итерационно решать именно ограниченную задачу,
обычный свободный градиентный шаг сам по себе не обеспечивал бы
допустимость. Потребовалось бы учитывать ограничение в алгоритме.
Результаты трёх аналитических постановок и безусловного численного
поиска нужно сравнивать с учётом их допустимых множеств.

\section{Сравнение методов и устный ответ на защите}

\question{Чем отличаются Ньютон и градиентный спуск}
Градиентный метод использует первые производные и выбирает направление
$-g_k$; его поведение сильно зависит от шага.
Ньютон использует также вторые производные и определяет поправку
через систему $H_ks=-g_k$.
Кривизна масштабирует разные направления по-разному.

В нашей задаче Ньютон достигает точного решения за один шаг,
градиентный метод --- приближённого решения за 23 или 36 шагов
при выбранных параметрах. Это не означает, что Ньютон всегда выгоднее
в любой задаче: формирование гессиана и решение системы могут быть
существенно дороже одного вычисления градиента.
Здесь размерность равна двум, а гессиан постоянен, поэтому пример
особенно благоприятен для Ньютона.

\question{Что нужно объяснить по графикам}
Линии уровня $J$ являются эллипсами с центром в безусловном минимуме.
Парабола $c=0$ задаёт границу, а закрашенная область --- допустимые
точки для неравенства. Безусловный минимум находится вне этой области;
ограниченный минимум расположен на границе.

На траекториях видно одно прямое перемещение Ньютона и последовательные
градиентные шаги. При $\gamma=0.09$ координаты переходят через
оптимальные значения с затухающей амплитудой.
Убывание нормы градиента показывает приближение к стационарности.

В полулогарифмическом масштабе линейная сходимость даёт приблизительно
прямые участки, потому что геометрическое убывание превращается
в линейное после логарифмирования.
Визуально плоский график $J$ ещё не означает нулевой градиент:
около минимума ошибка критерия имеет второй порядок по расстоянию.

\question{Какие числа и различия следует помнить}
\begin{center}
\small
\begin{tabular}{l|r}
\hline
Величина&Значение\\\hline
Безусловный минимум $z^*$&$(0,-0.25)^{\mathsf T}$\\
Безусловное значение $J^*$&$-10.375$\\
Ограниченный минимум $z^*$&$(-0.320399805,-0.196978643)^{\mathsf T}$\\
Ограниченное значение $J^*$&$-9.604836265$\\
Множитель $\lambda^*$&$1.638103816$\\
Условие остановки&$\norm{\nabla J}\leq10^{-6}$\\
Ньютон / шаг $0.05$ / шаг $0.09$&$1\ /\ 23\ /\ 36$ итераций\\\hline
\end{tabular}
\end{center}
Наиболее вероятные ошибки в объяснении:
\begin{itemize}
\item называть нулевой градиент достаточным условием минимума;
\item путать критерий $J$, лагранжиан $L$, градиент и гессиан;
\item считать, что шаг Ньютона нужно приравнять к нулю заранее;
\item объявлять любое активное ограничение имеющим $\lambda>0$;
\item применять достаточность ККТ для выпуклой задачи к нашей невыпуклой области;
\item называть постоянный шаг точным поиском минимума вдоль направления;
\item связывать колебания координат с обязательным ростом критерия;
\item считать, что увеличение шага всегда ускоряет сходимость.
\end{itemize}
Каждое из этих утверждений исправляется соответствующим условием
или доказательством, приведённым в ответах выше.

\question{Как связно представить работу за одну или две минуты}
«В работе я исследовал квадратичный критерий двух переменных.
Сначала нашёл безусловную стационарную точку из равенства градиента
нулю. Гессиан положительно определён, поэтому точка $(0,-0.25)$
является единственным глобальным минимумом, а значение критерия
равно $-10.375$.

Затем рассмотрел равенство $3x-u^2+1=0$. Условия Лагранжа и
подстановка $x=(u^2-1)/3$ приводят к одному кубическому уравнению.
После сведения к одной переменной вторая производная строго
положительна, поэтому минимум на параболе единственный.
Получена точка примерно $(-0.320400,-0.196979)$ с критерием $-9.604836$.

Для неравенства безусловный минимум оказался недопустимым,
а активная ветвь Куна--Таккера дала ту же граничную точку.
Глобальность доказана через положительно определённый гессиан
лагранжиана и нижнюю оценку критерия.

В численной части я искал безусловный минимум.
Ньютон нашёл его за один шаг, поскольку функция квадратична.
Градиентный метод с постоянным шагом $0.05$ сошёлся апериодически
за 23 итерации, а с шагом $0.09$ --- с затухающими колебаниями
за 36. Различие объясняется собственными значениями матрицы итерации.
Остановка проверялась по норме градиента, не большей $10^{-6}$».
\end{document}
